% ============================================================================
% TEORETICKÁ ČASŤ
% Podľa metodických pokynov FZ TnUAD - cca 30% práce
% Písané v 1. osobe množného čísla, minulý čas (autorský plurál)
% ============================================================================

\chapter{Súčasný stav riešenej problematiky}

\section{Reumatoidná artritída Reumatoidná artritída (RA) predstavuje závažné, chronické zápalové ochorenie autoimunitnej etiológie, ktoré sa vyznacuje systémovým charakterom a progresívnym priebe- ˇ hom. Jablkom sváru v patogenéze tohto ochorenia je imunitný systém, ktorý omylom identifikuje vlastné tkanivá ako cudzorodé, co vedie k deštrukcii, bolesti a strate funkcie. Hoci ˇ primárnym ciel’om patologického procesu sú periférne synoviálne k ́lby, najmä drobné sk ́lbenia rúk a nôh, RA sa v plnej miere prejavuje ako systémová choroba postihujúca aj mimok ́lbové štruktúry a vnútorné orgány. Tento aspekt ochorenia významne prispieva k zvýšenej morbidite a skráteniu d ́lžky života postihnutých jedincov. Ochorenie sa môže manifestovat’ v ktoromkol’vek veku, avšak demografické krivky jasne poukazujú na vrchol incidencie v produktívnom veku, typicky medzi 30. až a 50. rokom života, co má nesmierny celospolo ˇ censký ˇ a ekonomický dopad. Etiopatogenéza ochorenia je multifaktoriálna a hoci nie je do posledného detailu objasnená, predpokladá sa zložitá interakcia genetickej predispozície hostitel’a a spúšt’acích faktorov vonkajšieho prostredia. Priebeh choroby je vysoko individuálny – od miernych foriem s pomalou progresiou až po agresívne formy vedúce k rýchlej invalidizácii (Rovenský, 2001).
 Anatomickým substrátom, kde sa odohráva hlavná dráma autoimunitnej reakcie, je synoviálna membrána. Za normálnych okolností tenká vrstva buniek sa pri RA mení na hyperplastické, zapálené tkanivo prestúpené bunkami imunitného systému. Táto zmenená synovia, oznacovaná ako panus, sa správa ako lokálny tumor, ktorý invazívne prerastá do okolia, ˇ nahlodáva k ́lbovú chrupavku a následne aj subchondrálnu kost’. Výsledkom tohto procesu je nielen bolest’ a opuch, ale predovšetkým nezvratná deštrukcia k ́lbovej architektúry a strata hybnosti (Rovenský, 2001; Shah et al., 2024).
 1.1 Funkcná anatómia k ˇ  ́lbového systému Pochopenie anatomických a biomechanických princípov k ́lbového aparátu je nevyhnutným predpokladom pre správnu interpretáciu patologických zmien pri reumatoidnej artritíde. K ́lby l’udského tela možno rozdelit’ podl’a rôznych kritérií, pricom z pohl’adu reuma- ˇ tológie má primárny význam delenie podl’a prítomnosti k ́lbovej dutiny na spojenia nepravé (synartrózy) a pravé k ́lby (diartrózy). Práve pravé k ́lby obsahujúce synoviálnu membránu sú tercom autoimunitného útoku pri RA (Dylevský, ˇ 2009).
 1.1.1 Stavba synoviálneho k ́lbu Synoviálny k ́lb predstavuje zložitú anatomickú jednotku, ktorej základnými komponentmi sú k ́lbové povrchy pokryté hyalínnou chrupavkou, k ́lbové puzdro s jeho vrstvami a k ́lbová dutina vyplnená synoviálnou tekutinou. Každá z týchto štruktúr zohráva špecifickú úlohu v biomechanike pohybu a súcasne predstavuje potenciálny substrát pre patologické ˇ zmeny (Dylevský, 2009).
 K ́lbová chrupavka pokrýva artikulacné plochy kostí a zabezpe ˇ cuje hladký pohyb s ˇ minimálnym trením. Je tvorená prevažne extracelulárnou matrix bohatou na kolagén typu II a proteoglykány, ktoré viažu vodu a poskytujú chrupavke jej typické viskoelastické vlastnosti.
 Chondrocyty, jediné bunky prítomné v chrupavke, udržiavajú rovnováhu medzi syntézou a degradáciou matrix. Pri RA dochádza k narušeniu tejto rovnováhy v prospech degradácie, coˇ vedie k progresívnemu úbytku chrupavkového tkaniva.
 K ́lbové puzdro pozostáva z dvoch vrstiev. Vonkajšia fibrózna vrstva (membrana fibrosa) je tvorená hustým väzivom bohatým na kolagénové vlákna typu I, ktoré zabezpecujú ˇ mechanickú stabilitu k ́lbu. Vnútorná synoviálna vrstva (membrana synovialis) je pri RA miestom najintenzívnejších zápalových zmien. Za fyziologických podmienok je synoviálna membrána tenká, tvorená jednou až tromi vrstvami buniek – synoviocytov (Koláˇr, 2009).
 1.1.2 Synoviálna membrána a jej úloha Synoviálna membrána je metabolicky aktívne tkanivo, ktoré produkuje synoviálnu tekutinu – viskózny ultrafiltrátt krvnej plazmy obohatený o kyselinu hyalurónovú. Táto tekutina plní niekol’ko kl’úcových funkcií: lubrikáciu k ˇ  ́lbových povrchov, výživu avaskulárnej chrupavky difúziou živín a odstranovanie metabolických produktov z k ˇ  ́lbovej dutiny.
 V synoviálnej membráne rozlišujeme dva hlavné typy buniek. Synoviocyty typu A majú makrofágový pôvod a zabezpecujú fagocytózu debris a imunitný dohl’ad. Synoviocyty ˇ typu B sú fibroblastového charakteru a produkujú komponenty synoviálnej tekutiny, najmä kyselinu hyalurónovú a lubricín. Pri RA dochádza k masívnej infiltrácii synoviálnej membrány imunitnými bunkami a transformácii rezidentných synoviocytov do agresívneho fenotypu (Smith, 2015).
 1.1.3 Anatomické predilekcné miesta pri RA ˇ Reumatoidná artritída vykazuje charakteristickú predilekciu pre urcité k ˇ  ́lbové skupiny, co má diagnostický význam. Naj ˇ castejšie postihnuté sú metakarpofalangeálne (MCP) ˇ k ́lby, proximálne interfalangeálne (PIP) k ́lby rúk, zápästné k ́lby a metatarzofalangeálne (MTP) k ́lby nôh. Distálne interfalangeálne (DIP) k ́lby sú typicky ušetrené, co odlišuje RA od oste- ˇ oartrózy a psoriatickej artritídy (Olejárová, 2008).
 Táto selektívna distribúcia postihnutia nie je náhodná. Drobné k ́lby prstov obsahujú pomerne vel’kú synoviálnu membránu v pomere k vel’kosti k ́lbu, co poskytuje vä ˇ cšiu plo- ˇ chu pre autoimunitnú reakciu. Navyše tieto k ́lby sú vystavené opakovanému mechanickému stresu pri bežných denných cinnostiach, ˇ co môže potencovat’ zápalový proces (K. Pavelka; ˇ Rovenský et al., 2003).
 1.2 Epidemiológia Výskyt a rozšírenie RA sa líši podl’a diagnostických schém, pricom incidencia je ˇ vyššia u žien než u mužov, a to v 70 percentách prípadov. Prevalencia sa pohybuje okolo pol percenta až jedného percenta populácie, líši sa avšak podl’a regiónu a etnických skupín.
 Zatial’ co v Amerike je prevalencia vyššia, v Afrike je práve naopak nižšia. RA je charakte- ˇ ristická výskytom vo veku 20 až 55 rokov; s novými štúdiami sa potvrdilo, že vzniká už aj v neskoršom štádiu veku (Rovenský, 2001; Rovenský, 1998).
 Podl’a odhadov v roku 2020 trpelo RA celosvetovo približne 17,6 milióna obyvatel’ov, po zohl’adnení vekovej štruktúry populácie sa odhaduje, že celosvetovo postihuje reumatoidná artritída približne 0,21 % l’udí. V porovnaní s rokom 1990 došlo k 14,1% nárastu poctu prípadov RA. Predpokladá sa, že do 25 rokov po ˇ cet pacientov s RA stúpne na 31,7 mi- ˇ lióna. Epidemiologické štúdie v Cíne uvádzajú okolo 5 miliónov pacientov, ˇ co predstavuje ˇ prevalenciu 0,42 % populácie (Gao et al., 2020; Rovenský, 2001; Rovenský, 1998).
 U starších pacientov zaznamenávame odlišné epidemiologické charakteristiky. Gerontorevmatológia poukazuje na špecifiká RA v geriatrickej populácii, kde ochorenie casto ˇ nadobúda atypický klinický priebeh. Séropozitivita býva nižšia a diferenciálna diagnostika voci polymyalgii reumatike a osteoartróze je náro ˇ cnejšia (Rovenský, ˇ 2014).
 1.3 Etiológia a patofyziológia Súcasná veda nazerá na vznik reumatoidnej artritídy ako na viacstup ˇ nový proces, v ˇ ktorom dochádza k prelomeniu imunologickej tolerancie. Na zaciatku stojí geneticky pre- ˇ disponovaný jedinec (nosic ur ˇ citých alel HLA-DRB1), u ktorého environmentálny spúš- ˇ t’ac – napríklad faj ˇ cenie, chronická parodontitída (baktéria ˇ Porphyromonas gingivalis) alebo zmeny v crevnom mikrobióme – naštartuje produkciu autoprotilátok. Tento predklinický stav ˇ môže trvat’ roky predtým, než sa objaví prvá synovitída. Kl’úcovými biomarkerovými hrá ˇ cmi ˇ v tomto procese sú reumatoidný faktor (RF) a protilátky proti cyklickému citrulínovanému peptidu (anti-CCP), ktorých prítomnost’ casto predchádza klinickým symptómom (Shah et ˇ al., 2024; Scherer et al., 2020).
 1.3.1 Autoimunitné mechanizmy Autoimunita predstavuje stav, pri ktorom imunitný systém stráca schopnost’ rozlišovat’ medzi vlastnými a cudzorodými antigénmi. Pri fyziologických podmienkach existujú viaceré kontrolné mechanizmy – centrálna a periférna tolerancia – zabranujúce útoku na ˇ vlastné tkanivá. U pacientov s RA dochádza k zlyhaniu týchto mechanizmov na viacerých úrovniach (Buc, 2016).
 Genetická predispozícia zohráva kl’úcovú úlohu v rozvoji autoimunitného procesu. ˇ Asociácia s hlavným histokompatibilným komplexom (MHC), konkrétne s urcitými ale- ˇ lami HLA-DRB1 obsahujúcimi tzv. zdiel’aný epitop (shared epitope), je jedným z najsilnejších genetických rizikových faktorov. Tieto molekuly HLA prezentujú autoantigény Tlymfocytom, cím iniciujú patologickú imunitnú odpoved’ (Tak et al., ˇ 2011).
 Proces citrulinácie predstavuje ústredný mechanizmus v patogenéze RA. Ide o posttranslacnú modifikáciu proteínov, pri ktorej enzým peptidylarginín deimináza (PAD) kon- ˇ vertuje arginín na citrulín. Citrulínované proteíny sa stávajú neoantigénmi, voci ktorým or- ˇ ganizmus vytvára protilátky (ACPA). Tieto autoprotilátky sa môžu objavit’ roky pred manifestáciou klinických príznakov a ich prítomnost’ je silným prediktorom erozívneho priebehu ochorenia (Choy, 2012).
 1.3.2 Patofyziologická kaskáda V patofyziologickej kaskáde dominuje dysregulácia imunitného systému. Dochádza k patologickej aktivácii T-lymfocytov (najmä Th1 a Th17 subpopulácií), ktoré migrujú do k ́lbovej výstelky. Tu prostredníctvom medzibunkovej komunikácie stimulujú makrofágy, Blymfocyty a fibroblasty k nadprodukcii prozápalových mediátorov. Vzniká tzv. cytokínová búrka, v ktorej dominujú tumor nekrotizujúci faktor alfa (TNF-α), interleukín-1 (IL-1) a interleukín-6 (IL-6). Tieto molekuly sú zodpovedné nielen za lokálny zápal a bolest’, ale aj za systémové prejavy, ako je únava, horúcka ˇ ci tvorba reaktantov akútnej fázy v pe ˇ ceni ˇ (CRP) (Gao et al., 2020).
 1.3.3 Úloha synoviálnych fibroblastov Synoviálne fibroblasty, oznacované aj ako fibroblastom podobné synoviocyty (FLS), ˇ zohrávajú centrálnu úlohu v patogenéze RA. Na rozdiel od imunitných buniek, ktoré do k ́lbu infiltrujú zvonka, FLS sú rezidentné bunky synoviálnej membrány, ktoré pri RA podliehajú fenotypovej transformácii. Nadobúdajú agresívne, nádoru podobné vlastnosti charakterizované zvýšenou proliferáciou, zníženou apoptózou a schopnost’ou invadovat’ do chrupavky a kosti (Smith, 2015).
 Aktivované FLS produkujú široké spektrum deštruktívnych enzýmov, predovšetkým matrix metaloproteinázy (MMP-1, MMP-3, MMP-13), ktoré degradujú kolagén a d’alšie komponenty extracelulárnej matrix. Súcasne sekretujú chemokíny prit’ahujúce d’alšie zápa- ˇ lové bunky a rastové faktory podporujúce angiogenézu v zapálenom tkanive. Tento autokriný a parakriný komunikacný systém vytvára sebaposil ˇ nujúci cyklus chronického zápalu (Smith, ˇ 2015; Tak et al., 2011).
 1.3.4 Tvorba panusu a deštrukcia k ́lbu Chronická stimulácia vedie k transformácii synoviocytov, ktoré sa zacínajú nekon- ˇ trolovane množit’ a správat’ invazívne. Vytvára sa hypervaskularizované granulacné tka- ˇ nivo – panus. Bunky panusu produkujú metaloproteinázy a iné enzýmy, ktoré doslova ``trávia``k ́lbovú chrupavku. Súcasne dochádza k aktivácii osteoklastov, buniek odbúravajúcich ˇ kost’, co vedie k vzniku typických kostných erózií. Tento deštruktívny proces je nezvratný; ˇ raz znicenú chrupavku a kost’ organizmus nedokáže plnohodnotne regenerovat’, ˇ co pod ˇ ciar- ˇ kuje nutnost’ vcasnej a agresívnej lie ˇ cby (Rovenský, ˇ 2001; Isaacs et al., 2002).
 1.4 Klinický obraz Klinická manifestácia RA je povestná svojou heterogenitou. U väcšiny pacientov ˇ (približne 70 %) sa ochorenie vkráda do života plazivo, nenápadne, v priebehu niekol’kých týždnov ˇ ci mesiacov. Prvými príznakmi nemusia byt’ bolesti k ˇ  ́lbov, ale nešpecifické, prodromálne symptómy pripomínajúce virózu: nmerná únava, subfebrílie, nechutenstvo a mierny úbytok hmotnosti. Až následne sa pridružuje typická artikulárna symptomatológia. Dominantným znakom je ranná stuhnutost’ k ́lbov. Na rozdiel od krátkodobej stuhnutosti pri artróze, pri RA tento stav trvá dlhšie než hodinu, casto celé predpoludnie, a odznieva až po ˇ rozhýbaní. Stuhnutost’ je odrazom nahromadenia zápalového edému v k ́lbe pocas no ˇ cného ˇ pokoja (Nemec, ˇ 2023).
 Artritída sa klinicky prejavuje klasickou päticou príznakov zápalu: bolest’ou (dolor), opuchom (tumor), zvýšenou teplotou k ́lbu (calor) a poruchou funkcie (functio laesa); zacer- ˇ venanie (rubor) býva pri RA menej casté. Postihnutie je typicky polyartikulárne (zasahuje ˇ viac ako 3 k ́lbové oblasti) a symetrické (postihuje rovnaké k ́lby na oboch stranách tela).
 Predilekcnými lokalitami sú drobné k ˇ  ́lby ruky – metakarpofalangeálne (MCP) a proximálne interfalangeálne (PIP) k ́lby, ako aj zápästia a metatarzofalangeálne (MTP) k ́lby nohy. Distálne interfalangeálne k ́lby (DIP) bývajú typicky ušetrené, co je dôležitý diferencia ˇ cný znak ˇ oproti osteoartróze a psoriatickej artritíde (Nemec, ˇ 2023; Rovenský, 2001).
 U menšej casti pacientov môže mat’ choroba akútny, búrlivý nástup, ktorý pacienta ˇ doslova pripúta na lôžko v priebehu niekol’kých dní. Ešte zriedkavejší je palindromatický reumatizmus, charakterizovaný krátkymi, opakujúcimi sa záchvatmi artritídy, ktoré nezanechávajú trvalé následky, no môžu prejst’ do klasickej RA. Zvláštnu pozornost’ si vyžaduje RA u starších l’udí, ktorá sa casto prezentuje asymetrickým postihnutím vel’kých k ˇ  ́lbov, najmä ramenných pletencov (tzv. ``shoulder-likeönset), co môže viest’ k zámene s polymyal- ˇ giou reumaticou (Nemec, ˇ 2023).
 1.4.1 K ́lbové deformity a štrukturálne zmeny Dlhodobo neostatocne kontrolovaný zápal vedie k deštrukcii väzivového aparátu a ˇ k ́lbových puzdier, co má za následok vznik typických deformít. Tieto zmeny nie sú len koz- ˇ metickým defektom, ale vedú k t’ažkej funkcnej poruche úchopu a lokomo ˇ cnej schopnosti ˇ (Olejárová; kol., 2016).
 Ruka a zápästie: Artikulárny komplex ruky je zrkadlom RA. Synovitída zápästia vedie k deštrukcii karpálnych kosticiek a väzov, ˇ coˇ casto vyústi do volárnej subluxácie zápästia ˇ a instability. Castou komplikáciou je kompresia *nervus medianus* – syndróm karpálneho ˇ tunela. Na prstoch pozorujeme charakteristickú ulnárnu deviáciu (vybocenie prstov smerom ˇ k lakt’ovej kosti) v MCP k ́lboch. Dalšími typickými deformitami sú: ˇ • Deformita gombíkovej dierky (Boutonnière): Je charakterizovaná fixovanou flexiou v PIP k ́lbe a hyperextenziou v DIP k ́lbe. Vzniká pretrhnutím centrálneho pruhu extenzorovej šl’achy.
 • Deformita labutej šije (Swan neck): Opacná situácia – hyperextenzia v PIP k ˇ  ́lbe a flexia v DIP k ́lbe, spôsobená spazmom vnútorných svalov ruky a laxicitou volárnej platnicky. ˇ Noha: Postihnutie nôh je casté a bolestivé. Dochádza k poklesu prie ˇ cnej klenby, ˇ vzniku hallux valgus (vybocený palec) a kladivkových prstov (digiti mallei). Pod hlavi ˇ ckami ˇ metatarzov sa tvoria bolestivé otlaky, ktoré pacientovi st’ažujú chôdzu a vyžadujú ortopedickú obuv.
 Vel’ké nosné k ́lby: Koleno je postihnuté casto, s tvorbou rozsiahlych výpotkov. Hy- ˇ pertrofická synovia alebo tekutina sa môže pretlacit’ do podkolennej jamky a vytvorit’ Bake- ˇ rovu cystu, ktorej prasknutie imituje trombózu žíl lýtka. Bedrový k ́lb býva postihnutý menej casto, ale koxitída má tendenciu k rýchlej progresii a vyžaduje ˇ casto totálnu endoprotézu. ˇ Chrbtica: Kritickou oblast’ou je krcná chrbtica. Zápal v okolí atlantoaxiálneho sk ˇ  ́lbenia (C1-C2) môže viest’ k uvol’neniu väzov a instabilite. Tento stav je život ohrozujúci, pretože pri neopatrnom pohybe môže dôjst’ k útlaku miechy. Každý pacient s RA podstupujúci operáciu v celkovej anestézii musí mat’ RTG krcnej chrbtice na vylú ˇ cenie tejto komplikácie. ˇ Obr. 1.1: Schematické znázornenie typických deformít ruky pri pokrocilej reumatoidnej ar- ˇ tritíde. Zobrazená je ulnárna deviácia v MCP k ́lboch a deformity prstov (Netter, 2014) 1.4.2 Mimok ́lbové (extraartikulárne) prejavy Reumatoidná artritída je systémové ochorenie par excellence. Prítomnost’ mimok ́lbových prejavov signalizuje t’ažšiu formu ochorenia a horšiu prognózu (Rovenský, 2001).
 Reumatoidné uzly: Sú najcastejším extraartikulárnym prejavom, vyskytujúcim sa u ˇ 20-30 % séropozitívnych pacientov. Sú to tuhé, nebolestivé podkožné útvary, ktoré sa tvoria najmä v miestach mechanického tlaku – nad lakt’ovým výbežkom (olecranon), na chrbte ruky ci na achilovke. ˇ Ocné postihnutie: ˇ Najbežnejšia je keratokonjunktivitída sicca (syndróm suchého oka) ako súcast’ sekundárneho Sjögrenovho syndrómu. Závažnejšou komplikáciou je sk- ˇ leritída alebo episkleritída, ktorá sa prejavuje bolestivým zacervenaním oka a môže viest’ až ˇ k poškodeniu zraku.
 Pl’úcne postihnutie: Pl’úca sú castým ter ˇ com. Môže íst’ o pleuritídu (zápal pohrud- ˇ nice) s výpotkom, alebo o závažnejšiu intersticiálnu pl’úcnu fibrózu, ktorá postupne znižuje dýchaciu kapacitu pl’úc.
 Vaskulitída: Reumatoidná vaskulitída je zápal ciev, ktorý môže viest’ k tvorbe kožných vredov, nekróz na prstoch alebo k postihnutiu nervov (mononeuritis multiplex). Je to vážna komplikácia vyžadujúca agresívnu liecbu. ˇ 1.5 Diagnostika Moderná reumatológia kladie dôraz na vcasnú diagnostiku, tzv. ökno príležitosti``(window ˇ of opportunity). Ciel’om je zachytit’ ochorenie v jeho pociatkoch, kedy je terapeutická inter- ˇ vencia najúcinnejšia a dokáže zabránit’ trvalému poškodeniu k ˇ  ́lbov. Diagnostický proces je komplexný a vychádza zo syntézy klinického obrazu, laboratórnych výsledkov a zobrazovacích metód. Zlatým štandardom sú v súcasnosti ˇ klasifikacné kritériá ACR/EULAR z roku ˇ 2010, ktoré boli vyvinuté práve pre záchyt vcasných štádií artritídy (Aletaha et al., ˇ 2010).
 Bodovací systém kritérií hodnotí štyri domény, pricom na potvrdenie diagnózy ``jed- ˇ noznacnej RA``je potrebný sú ˇ cet ˇ 6 a viac bodov z maximálnych 10: 1. Postihnutie k ́lbov (0 – 5 bodov): Dôraz sa kladie na pocet a typ postihnutých k ˇ  ́lbov.
 Najvyššie bodové ohodnotenie (5 bodov) získava pacient s postihnutím viac ako 10 k ́lbov, pricom aspo ˇ n jeden z nich musí byt’ malý k ˇ  ́lb.
 2. Sérológia (0 – 3 body): Prítomnost’ reumatoidných faktorov (RF) a protilátok proti citrulínovaným peptidom (ACPA/anti-CCP). Vysoké titre týchto protilátok sú vysoko špecifické pre RA a prognosticky nepriaznivé.
 3. Reaktanty akútnej fázy (0 – 1 bod): Laboratórny dôkaz systémového zápalu – zvýšená sedimentácia erytrocytov (FW) alebo C-reaktívny proteín (CRP).
 4. Trvanie symptómov (0 – 1 bod): Pretrvávanie t’ažkostí dlhšie ako 6 týždnov odlišuje ˇ chronickú artritídu od prechodných vírusových artralgií.
 Okrem klasického röntgenového vyšetrenia (RTG), ktoré v skorom štádiu môže byt’ negatívne, sa dnes do popredia dostávajú senzitívnejšie metódy. Ultrasonografia (USG) k ́lbov s Power Doppler módom dokáže zobrazit’ synoviálny zápal a prekrvenie skôr, než dôjde k poškodeniu kosti. Magnetická rezonancia (MRI) je najcitlivejšou metódou na detekciu kostného edému, ktorý je predzvest’ou budúcich erózií (Smolen et al., 2018).
 1.6 Komplikácie reumatoidnej artritídy Je nevyhnutné vnímat’ RA nielen ako izolovanú artropatiu, ale ako komplexnú systémovú chorobu, ktorej zápalový potenciál indukuje rad závažných komorbidít a komplikácií, majúcich zásadný vplyv na morbiditu a predcasnú mortalitu postihnutých jedincov. Tieto ˇ nežiaduce stavy môžu pramenit’ priamo z aktivity základného ochorenia alebo môžu byt’ iatrogénnym následkom dlhodobej imunosupresívnej liecby (Smolen et al., ˇ 2018).
 1.6.1 Kardiovaskulárne komplikácie Epidemiologické štúdie konzistentne potvrdzujú, že populácia s RA celí markantne ˇ vyššiemu riziku kardiovaskulárnych príhod (približne 1,5 až 2-násobne) v porovnaní so zdravou populáciou. Práve kardiovaskulárne zlyhanie je hlavnou prícinou skrátenia d ˇ  ́lžky života u týchto pacientov. Pretrvávajúci systémový zápal poškodzuje endotelovú výstelku ciev a akceleruje aterosklerotické procesy, co vyúst’uje do zvýšenej incidencie infarktu myokardu, ˇ cievnych mozgových príhod a kongestívneho srdcového zlyhania (Smolen et al., 2018).
 Mechanizmy kardiovaskulárneho poškodenia Patofyziologický vzt’ah medzi reumatoidnou artritídou a kardiovaskulárnym ochorením je komplexný a zah ́rna viaceré mechanizmy. Chronický systémový zápal prispieva k ˇ endotelovej dysfunkcii – pociato ˇ cnému štádiu aterosklerózy. Prozápalové cytokíny (TNF- ˇ α, IL-6) aktivujú endotelové bunky, cím podporujú adhéziu monocytov a ich transmigráciu do ˇ cievnej steny. Výsledkom je formácia aterómových plátov, ktoré sú u pacientov s RA casto ˇ menej stabilné a náchylnejšie na ruptúru (Freedman et al., 2014).
 Rizikový profil je navyše zhoršený prítomnost’ou konvencných faktorov, ako sú hy- ˇ pertenzia a dyslipidémia, ktoré môžu byt’ potencované terapiou kortikoidmi a nesteroidnými antireumatikami (NSA). Z klinického hl’adiska je významná aj perikarditída, ktorá hoci casto ˇ prebieha asymptomaticky, môže byt’ detegovaná echokardiograficky. Striktný manažment krvného tlaku a lipidov je preto imperatívom (Smolen et al., 2018; Freedman et al., 2014).
 Preventívne stratégie Prevencia kardiovaskulárnych príhod u pacientov s RA vyžaduje komplexný prístup.
 Na prvom mieste je efektívna kontrola zápalovej aktivity základného ochorenia – pacienti v remisii majú významne nižšie kardiovaskulárne riziko. Dalej je potrebný striktný manaž- ˇ ment tradicných rizikových faktorov vrátane lie ˇ cby hypertenzie, dyslipidémie a diabetu. Od- ˇ porúca sa používat’ modifikované skórovacie systémy (SCORE) s korek ˇ cným faktorom 1,5 ˇ pre pacientov s RA. Významná je aj úloha pravidelnej fyzickej aktivity, ktorá nielen zlepšuje kardiovaskulárnu kondíciu, ale má aj protizápalové úcinky (Freedman et al., ˇ 2014).
 1.6.2 Osteoporóza a fraktúry Demineralizácia kostného tkaniva, osteoporóza, predstavuje jednu z najcastejších ex- ˇ traartikulárnych komplikácií s komplexnou etiológiou. V patogenéze rozlišujeme lokálnu periartikulárnu osteoporózu, ktorá je priamym následkom pôsobenia prozápalových cytokínov v okolí postihnutého k ́lbu, a generalizovanú systémovú osteoporózu. Na jej vzniku sa podiel’a synergický efekt chronického zápalu, imobilizácie pacienta (zníženie mechanickej zát’aže kosti pre bolest’) a najmä dlhodobej kortikoterapie, ktorá inhibuje osteoblastickú aktivitu a redukuje intestinálnu resorpciu vápnika. Dôsledkom je zvýšená fragilita skeletu a vysoké riziko patologických fraktúr, predovšetkým stavcov a proximálneho femuru, co vedie ˇ k d’alšiemu zhoršeniu funkcného stavu (Karel Pavelka et al., ˇ 2018).
 1.6.3 Infekcné komplikácie ˇ Dysfunkcia imunitného dohl’adu, ktorá je inherentnou súcast’ou patofyziológie RA, ˇ v kombinácii s úcinkami potentných imunosupresív (vrátane biologík), robí pacientov zra- ˇ nitel’nými voci infek ˇ cným agens. Infekcie respira ˇ cného traktu (pneumónie), urogenitálneho ˇ systému a kože majú u týchto pacientov tendenciu k t’ažšiemu priebehu a castejším kom- ˇ plikáciám. Pri liecbe inhibítormi TNF- ˇ α existuje špecifické riziko reaktivácie latentnej tuberkulózy alebo chronickej hepatitídy B, co vyžaduje prísny skríning pred zahájením lie ˇ cby. ˇ Zvýšená je aj incidencia herpes zoster. Vakcinácia proti chrípke a pneumokokom je preto kl’úcovou preventívnou stratégiou (Macejová, ˇ 2019).
 1.6.4 Pl’úcne a hematologické komplikácie Pl’úcne prejavy sú castou viscerálnou manifestáciou RA. Môžu mat’ formu intersti- ˇ ciálnej pl’úcnej choroby (ILD), ktorá progreduje do nezvratnej fibrózy, pleurálnych výpotkov alebo reumatoidných uzlov v pl’úcnom parenchýme (v kombinácii s pneumokoniózou známe ako Caplanov syndróm). Klinickým korulátom je progredujúca dyspnoe a suchý, dráždivý kašel’. V krvnom obraze dominuje normocytárna anémia chronických chorôb, mechanizmom ktorej je hepcidínom sprostredkovaná blokáda metabolizmu železa. Zriedkavou, no život ohrozujúcou entitou je Feltyho syndróm (triáda: RA, splenomegália, granulocytopénia), ktorý extrémne zvyšuje náchylnost’ na fulminantné bakteriálne infekcie (Rovenský, 2001).
 Intersticiálne ochorenie pl’úc pri RA Intersticiálne ochorenie pl’úc asociované s reumatoidnou artritídou (RA-ILD) predstavuje závažnú extraartikulárnu manifestáciu, ktorá významne zhoršuje prognózu pacientov.
 Prevalencia subklinického postihnutia pl’úc detegovaného pomocou vysokorozlišovacieho CT (HRCT) dosahuje až 60 percent, pricom klinicky manifestná forma sa vyskytuje u 10 až ˇ 30 percent chorých. Rizikové faktory zah ́rnajú mužské pohlavie, vyšší vek, faj ˇ cenie, vysoké ˇ titre ACPA protilátok a prítomnost’ reumatoidných uzlov (Kadura et al., 2021).
 Patogenéza RA-ILD zah ́rna komplexnú interakciu genetických faktorov, autoimunit- ˇ ných mechanizmov a environmentálnych vplyvov. Cytokíny produkované v rámci systémového zápalu, najmä TNF-alfa, IL-6 a transformujúci rastový faktor beta (TGF-β), stimulujú proliferáciu fibroblastov a nadmernú depozíciu kolagénu v pl’úcnom interstíciu. Histopatologicky prevláda obraz obvyklej intersticiálnej pneumónie (UIP) alebo nešpecifickej intersticiálnej pneumónie (NSIP) (Kadura et al., 2021; Klener, 2011).
 Manažment RA-ILD vyžaduje multidisciplinárny prístup zah ́rnajúci reumatológa a ˇ pneumológa. Terapeutické možnosti sú limitované; niektoré DMARDs (napr. metotrexát) môžu potenciálne zhoršit’ pl’úcne postihnutie, zatial’ co iné (rituximab, abatacept) sa javia ˇ ako bezpecnejšie. V pokro ˇ cilých štádiách prichádza do úvahy antifibrotická lie ˇ cba pirfeni- ˇ dónom alebo nintedanibom (Kadura et al., 2021).
 1.6.5 Metabolické poruchy pri RA Vzt’ah medzi reumatoidnou artritídou a metabolickými poruchami je bidirecionálny a komplexný. Chronický zápal pri RA negatívne ovplyvnuje metabolické parametre, zatial’ ˇ coˇ metabolický syndróm zhoršuje zápalovú aktivitu a komplikuje liecbu základného ochorenia ˇ (Kašperová et al., 2021).
 Inzulínová rezistencia a diabetes mellitus Pacienti s RA majú zvýšené riziko rozvoja diabetes mellitus 2. typu v porovnaní s bežnou populáciou. Prozápalové cytokíny, predovšetkým TNF-α a IL-6, interferujú so signálnou dráhou inzulínu na úrovni inzulínového receptorového substrátu (IRS-1), co vedie ˇ k rozvoju inzulínovej rezistencie. Situáciu d’alej komplikuje dlhodobá kortikoterapia, ktorá priamo zvyšuje glykémiu a podporuje viscerálnu obezitu (Kašperová et al., 2021).
 Dyslipidémia a ateroskleróza Lipidový profil u pacientov s aktívnou RA vykazuje charakteristické zmeny oznacované ako ``lipidový paradox``. Napriek nízkym hodnotám celkového cholesterolu a LDL- ˇ cholesterolu (ktoré sú konzumované zápalovým procesom) je kardiovaskulárne riziko zvýšené. Prícinou sú kvalitatívne zmeny lipoproteínov – oxidácia LDL ˇ castíc a znížená antiatero- ˇ génna aktivita HDL cholesterolu. Efektívna kontrola zápalovej aktivity pomocou DMARDs vedie k paradoxnému zvýšeniu hladín lipidov, co však reflektuje zlepšenie metabolického ˇ stavu (Freedman et al., 2014; Kašperová et al., 2021).
 Reumatická kachexia Osobitným metabolickým fenoménom pri RA je reumatická kachexia, charakterizovaná stratou svalovej hmoty pri zachovanej alebo dokonca zvýšenej tukovej hmote. Na rozdiel od klasickej kachexie pri malígnych ochoreniach, telesná hmotnost’ môže zostat’ nezmenená, co maskuje závažný úbytok svalstva. Patofyziologicky sa uplat ˇ nuje katabolický ˇ efekt prozápalových cytokínov, znížená fyzická aktivita pre bolest’ a artritídu a niekedy aj nedostatocný príjem bielkovín. Dôsledkom je svalová slabost’, zhoršená funk ˇ cná kapacita a ˇ zvýšené riziko pádov (Klener, 2011).
 Dopad reumatoidnej artritídy na kvalitu života (QoL) je devastacný a multodimen- ˇ zionálny, zasahujúci d’aleko za hranice fyzickej disability. WHO definuje zdravie ako stav kompletnej bio-psycho-sociálnej pohody, co je ideál, ktorý je u pacientov s RA ˇ casto nedo- ˇ siahnutel’ný (Krchnavá et al., ˇ 2018).
 1.6.6 Fyzická a funkcná doména ˇ Kl’úcovými faktormi, ktoré negatívne determinujú QoL, sú perzistujúca bolest’, pato- ˇ logická únava a ranná stuhnutost’. Bolest’, prítomná v pokoji aj pri aktivite, narúša spánkovú architektúru a vedie k syndrómu chronického vycerpania (fatigue), ktorý pacienti ˇ casto hod- ˇ notia ako viac limitujúci než bolest’ samotnú. Ranná stuhnutost’, trvajúca nezriedka hodiny, blokuje ranné rutinné cinnosti. Progresívna strata k ˇ  ́lbových rozsahov a svalovej sily znemožnuje sebaobsluhu (hygiena, obliekanie), ˇ co vedie k závislosti na pomoci druhej osoby, strate ˇ autonómie a pocitom bezmocnosti. Vzniká tak circulus vitiosus: bolest’ – inaktivita – atrofia – zhoršenie funkcie (Krchnavá et al., ˇ 2018).
 1.6.7 Psychologická doména Chronický, celoživotný charakter choroby s nepredvídatel’ným striedaním relapsov a remisií, spolu s hrozbou trvalého postihnutia, generuje u pacientov chronickú úzkost’ a depresívne poruchy. Prevalencia depresie sa odhaduje na 15 – 40 %, pricom je známe, že ˇ depresia znižuje prah bolesti a zhoršuje adherenciu k liecbe. Významným faktorom je aj ˇ alterácia telesného obrazu (body image) v dôsledku viditel’ných deformít, co vedie k so- ˇ ciálnemu stiahnutiu. Pacienti casto zažívajú fenomén ``nau ˇ cenej bezmocnostiä stratu identity ˇ (Olejárová; kol., 2016).
 1.6.8 Sociálna a ekonomická doména RA postihuje jedincov na vrchole produktívneho veku, co vedie k vysokej miere in- ˇ validizácie a strate pracovnej schopnosti. Absencia príjmu prináša ekonomickú instabilitu.
 Znížená mobilita vedie k sociálnej izolácii a zanechaniu zál’ub. Ochorenie mení dynamiku rodinných vzt’ahov, ked’ sú partneri nútení prevziat’ rolu opatrovatel’ov. Úlohou fyzioterapie je prostredníctvom edukácie a tréningu kompenzacných mechanizmov maximalizovat’ ˇ funkcnú nezávislost’ a tým zlepšit’ QoL (Krch ˇ navá et al., ˇ 2018; Poliaková et al., 2019).
 1.7 Diferenciálna diagnostika Proces diferenciálnej diagnostiky pri RA je intelektuálnou výzvou, ktorá vyžaduje rigorózne vylúcenie iných nozologických jednotiek s prekrývajúcou sa symptomatológiou. ˇ Správna a vcasná diferenciácia je sine qua non pre vol’bu adekvátnej terapie, nakol’ko lie- ˇ cebné prístupy sa pri jednotlivých ochoreniach diametrálne líšia (Karel Pavelka et al., ˇ 2018).
 1.7.1 Osteoartróza (OA) Osteoartróza je najcastejším ïmitátorom``RA, predovšetkým u geriatrickej populácie. ˇ Na rozdiel od RA, ktorá je primárne zápalovým ochorením, je OA degeneratívny proces spojený s opotrebovaním chrupavky.
 • Charakter bolesti: Pri OA je bolest’ mechanická – zhoršuje sa námahou a v pokoji ustupuje. Pri RA je bolest’ zápalová, prítomná aj v kl’ude.
 • Stuhnutost’: Pri OA je krátka, ``štartovacia``(do 30 minút). Pri RA je dlhodobá.
 • Lokalizácia: OA postihuje typicky DIP k ́lby (Heberdenove uzly), PIP k ́lby (Bouchardove uzly) a k ́lb palca ruky (rhizartróza). RA postihuje MCP k ́lby a zápästie.
 • Laboratórium: Pri OA sú zápalové markery (CRP, FW) v norme (Olejárová; kol., 2016).
 1.7.2 Séronegatívne spondyloartritídy Ide o skupinu zápalových ochorení, pre ktoré je typická neprítomnost’ reumatoidného faktora a väzba na antigén HLA-B27.
 1. Psoriatická artritída (PsA): Môže byt’ klinicky vel’mi podobná RA. Typickými diferenciacnými znakmi sú prítomnost’ psoriázy na koži (hoci artritída môže predchádzat’ ˇ kožným prejavom), postihnutie DIP k ́lbov, daktylitída (párkovitý opuch celého prsta) a zmeny na nechtoch (jamky, onycholýza). Na RTG vidíme súcasne erózie aj novotvorbu ˇ kosti (Macejová, 2019).
 2. Ankylózujúca spondylitída (Bechterevova choroba): Dominuje postihnutie axiálneho skeletu – zápalová bolest’ chrbta a sakroiliitída. Periférna artritída je zvycajne ˇ asymetrická a týka sa vel’kých k ́lbov dolných koncatín (Gravaldi et al., ˇ 2022).
 3. Reaktívna artritída: Vzniká v casovej súvislosti s prekonanou urogenitálnou alebo ˇ crevnou infekciou. Klasická Reiterova triáda zah ˇ  ́rna artritídu, uretritídu a konjunktivi- ˇ tídu (Scherer et al., 2020).
 1.7.3 Kryštálové artropatie (Dna) Dnavá artritída je metabolické ochorenie spôsobené depozíciou kryštálov nátriumurátu v k ́lboch.
 • Priebeh: Typický je akútny, perakútny dnavý záchvat s extrémne silnou bolest’ou, opuchom a zacervenaním, ktorý vzniká náhle, ˇ casto v noci (``podagra``– postihnutie ˇ MTP k ́lbu palca nohy).
 • Chronická forma: Tofózna dna môže imitovat’ RA prítomnost’ou uzlov (tofov) a erózií, ale rozlíšenie umožní punkcia synoviálnej tekutiny a dôkaz kryštálov (Rovenský, 2001).
 1.7.4 Systémový lupus erythematosus (SLE) SLE je multisystémové autoimunitné ochorenie.
 • K ́lby: Artritída pri SLE býva polyartikulárna a symetrická ako pri RA, ale je typicky neerozívna (Jaccoudova artropatia) – nenicí kost’, hoci spôsobuje deformity. ˇ • Ostatné príznaky: Fotosenzitivita, motýl’ovitý exantém na tvári, postihnutie obliciek ˇ (nefritída), serozitídy. Typická je prítomnost’ anti-dsDNA a anti-Sm protilátok (Rovenský et al., 2015).
 1.8 Terapia reumatoidnej artritídy Manažment RA prešiel v posledných dvoch dekádach revolucnými zmenami. Sú- ˇ casná stratégia je založená na koncepte ˇ ``Treat to Target`` (liecba k ciel’u), ktorého im- ˇ peratívom je dosiahnutie klinickej remisie alebo aspon stavu nízkej aktivity choroby v ˇ coˇ najkratšom case. Tento prístup si vyžaduje proaktívnu úpravu lie ˇ cby a úzku spoluprácu re- ˇ umatológa, pacienta a rehabilitacného tímu (Smolen et al., ˇ 2018).
 1.8.1 Farmakologická liecba ˇ Farmakoterapia zostáva pilierom liecby a delí sa na symptomatickú, ktorá len tlmí ˇ prejavy, a chorobu modifikujúcu, ktorá mení priebeh ochorenia.
 Symptomatická liecba (NSA a Kortikoidy) ˇ Tieto lieky poskytujú pacientovi rýchlu úl’avu, ale neriešia podstatu – nezastavujú k ́lbovú deštrukciu.
 • Nesteroidné antiflogistiká (NSA): Efektívne tlmia bolest’ a rannú stuhnutost’ inhibíciou enzýmov cyklooxygenázy. Ich dlhodobé užívanie je však limitované rizikom gastropatie (vredy), kardiovaskulárnych príhod a renálneho poškodenia.
 • Glukokortikoidy: Majú silný a rýchly protizápalový efekt. Využívajú sa najmä ako premost’ujúca liecba (bridging therapy) vo v ˇ casných štádiách alebo pri vzplanutiach. ˇ Pre riziko závažných nežiaducich úcinkov (osteoporóza, diabetes, hypertenzia, kata- ˇ rakta) je snahou podávat’ co najnižšiu ú ˇ cinnú dávku ˇ co najkratší ˇ cas (Olejárová; kol., ˇ 2016).
 Chorobu modifikujúce lieky (DMARDs) DMARDs (Disease Modifying Antirheumatic Drugs) zasahujú priamo do patogenetických mechanizmov a spomal’ujú rádiografickú progresiu.
 1. Konvencné syntetické DMARDs (csDMARDs): Metotrexát ˇ je ``zlatým štandardomä liekom prvej vol’by. Pôsobí ako antagonista kyseliny listovej. K d’alším liekom patria leflunomid, sulfasalazín a antimalariká (hydroxychlorochín).
 2. Biologické DMARDs (bDMARDs): Sú to geneticky inžinierované proteíny, ktoré cielene blokujú špecifické cytokíny alebo bunky. Sú indikované pri zlyhaní konvencnej ˇ liecby. Patrí sem: ˇ • anti-TNF liecba (Adalimumab, Infliximab), ˇ • blokátory IL-6 (Tocilizumab), • B-bunková deplécia (Rituximab), • modulátory kostimulácie T-buniek (Abatacept).
 3. Cielené syntetické DMARDs (tsDMARDs): Ide o tzv. malé molekuly (inhibítory JAK-kináz – tofacitinib, baricitinib), ktoré sa podávajú perorálne a blokujú vnútrobunkový prenos signálu. Sú rovnako úcinné ako biologiká (Scherer et al., ˇ 2020).
 1.9 Komplexná rehabilitácia Rehabilitácia je neoddelitel’nou súcast’ou terapeutického plánu. Zatial’ ˇ co farmako- ˇ terapia kontroluje zápal, rehabilitácia rieši funkcné dopady ochorenia. Jej ciel’om je udržat’ ˇ rozsah pohybu, svalovú silu, zmiernit’ bolest’ a naucit’ pacienta žit’ s chorobou tak, aby bol ˇ co najdlhšie sebesta ˇ cný (Kolá ˇ ˇr; kol., 2012).
 1.9.1 Kinezioterapia a liecebná telesná výchova ˇ Pohybová liecba (LTV) musí byt’ prísne individuálna a rešpektovat’ aktuálnu ak- ˇ tivitu ochorenia (akútne vs. chronické štádium). Nedávne systematické prehl’ady a metaanalýzy potvrdzujú, že pravidelné cvicenie je bezpe ˇ cné a má signifikantný pozitívny vplyv ˇ na funkcnú kapacitu a aeróbnu zdatnost’ bez zhoršenia aktivity ochorenia alebo poškodenia ˇ k ́lbov (Zhang et al., 2025; Metsios et al., 2020).
 Cvicenie v akútnom štádiu ˇ V case exacerbácie (vzplanutia) zápalu je prioritou ochrana ˇ k ́lbu (Koláˇr; kol., 2012).
 • Polohovanie: K ́lby sa polohujú vo funkcne výhodných postaveniach na prevenciu ˇ flekcných kontraktúr. ˇ • Pasívne pohyby: Vykonávajú sa opatrne do prahu bolesti na udržanie k ́lbovej vôl’e a výživy chrupavky.
 • Izometrické cvicenie: ˇ Svalové kontrakcie bez zmeny d ́lžky svalu sú kl’úcové na pre- ˇ venciu rýchlej svalovej atrofie, ktorá sprevádza zápal.
 (Koláˇr; kol., 2012; Navrátil et al., 2019) Cvicenie v chronickom štádiu a remisii ˇ V tomto období je ciel’om obnova funkcie a kondície (Metsios et al., 2020).
 1. Aeróbny tréning: Chôdza, cyklistika, plávanie. Odporúca sa stredná intenzita (60 – ˇ 70 % maximálnej tepovej frekvencie) 30 minút denne. Aeróbne cvicenie znižuje kar- ˇ diovaskulárne riziko, ktoré je u pacientov s RA zvýšené, a redukuje únavu (Metsios et al., 2020).
 2. Silový tréning: Cvicenie proti odporu (s pružnými t’ahmi Thera-Band, l’ahkými ˇ cin- ˇ kami) je nevyhnutné na boj proti reumatickej kachexii (úbytku svalovej hmoty). Štúdie ukazujú, že dynamický silový tréning je bezpecný a vedie k zlepšeniu svalovej sily a ˇ funkcnej schopnosti (Zhang et al., ˇ 2025).
 3. Cvicenie ruky: ˇ Špecifické posilnovacie a mobiliza ˇ cné cvi ˇ cenia pre drobné k ˇ  ́lby ruky (program SARAH) zlepšujú silu úchopu a manuálnu zrucnost’ efektívnejšie ako bežná ˇ starostlivost’ (Williams et al., 2015).
 High-intensity interval training Vysokointenzívny intervalový tréning (HIIT) predstavuje inovatívny prístup, ktorý bol dlho považovaný za nevhodný pre pacientov so zápalovými artropatiami. Nedávne randomizované kontrolované štúdie však priniesli prekvapivé zistenia. Štúdia ExeHeart preukázala, že HIIT v primárnej fyzioterapeutickej starostlivosti je bezpecný a efektívny u pacientov so zápalovou artritídou vrátane RA. Výsledky nazna- ˇ cujú zlepšenie aeróbnej kapacity a kardiovaskulárnych parametrov bez zhoršenia zápalovej ˇ aktivity ochorenia (Nordén et al., 2024).
 Kl’úcom k úspešnej implementácii HIIT u pacientov s RA je individuálna adaptá- ˇ cia protokolu. Intervaly vysokej intenzity by mali byt’ kratšie a odpocinkové fázy dlhšie v ˇ porovnaní so zdravou populáciou. Cvicenie by malo prebiehat’ na nebolestivých k ˇ  ́lboch s preferenciou aktivít s nízkou nárazovou zát’ažou, ako je cyklistika alebo plávanie (Nordén et al., 2024; Metsios et al., 2020).
 Podpora zmeny životného štýlu Fyzioterapeutom vedené intervencie zamerané na zmenu správania (behavior change interventions) sa ukazujú ako sl’ubný prístup na zvýšenie dlhodobej adherencie k pohybovej aktivite. Štúdia PIPPRA preukázala realizovatel’nost’ takejto intervencie u pacientov s RA. Kombinácia edukácie, stanovenia individuálnych ciel’ov a pravidelného follow-upu vedie k udržatel’nému zvýšeniu fyzickej aktivity (Larkin et al., 2024).
 1.9.2 Efektívnost’ fyzioterapie pri RA Systematické prehl’ady a meta-analýzy poskytujú vysokú úroven dôkazov pre ú ˇ cin- ˇ nost’ fyzioterapie v manažmente RA. Najnovší systematický prehl’ad analyzujúci randomizované kontrolované štúdie dospel k záveru, že fyzioterapeutická intervencia vedie k signifikantnému zmierneniu bolesti a zlepšeniu funkcných parametrov u pacientov s RA (El ˇ Ammare et al., 2023).
 Špecificky pre rehabilitáciu pri reumatických ochoreniach existujú štandardizované protokoly zah ́rnajúce kombináciu kinezioterapie, fyzikálnych procedúr a edukácie. Kom- ˇ plexný rehabilitacný program by mal byt’ zostavený multidisciplinárnym tímom a prispôso- ˇ bený individuálnym potrebám pacienta, fáze ochorenia a prítomnosti komorbidít (Jarošová et al., 2018).
 1.9.3 Fyzikálna terapia Fyzikálna terapia predstavuje integrálnu súcast’ rehabilita ˇ cného programu pri RA. ˇ Jej základným princípom je využitie fyzikálnych podnetov – tepla, chladu, elektrických impulzov, svetla a mechanickej energie – na dosiahnutie terapeutických úcinkov v ciel’ových ˇ tkanivách. Tieto procedúry primárne slúžia ako príprava pred kinezioterapiou, pricom ich ˇ analgezujúci a myorelaxacný efekt umož ˇ nuje pacientovi efektívnejšie cvi ˇ cit’ (Kraft, ˇ 2018; Kˇrížová, 2014).
 Termoterapia Aplikácia tepla (hypertermoterapia) nachádza uplatnenie predovšetkým v chronickom štádiu ochorenia. Teplo spôsobuje vazodilatáciu, zlepšuje perfúziu tkanív a uvol’nuje ˇ svalový spazmus. Najcastejšie využívané modality zah ˇ  ́rnajú parafínové zábaly, ktoré sú ob- ˇ zvlášt’ vhodné pre drobné k ́lby rúk, a solux lamp. V akútnom štádiu zápalu je naopak indikovaná kryoterapia (chlad), ktorá tlmí zápalovú reakciu a znižuje opuch (Vavrová et al., 2016).
 Elektroliecba ˇ Elektroterapeutické modality majú v liecbe RA dlhodobú tradíciu. TENS (transku- ˇ tánna elektrická nervová stimulácia) poskytuje analgéziu prostredníctvom vrátkového mechanizmu tlmenia bolesti. Interferencné prúdy (IFP) kombinujú výhody stredofrekven ˇ cných ˇ a nízkofrekvencných prúdov a sú suverénnou modalitou pre hlboko uložené k ˇ  ́lby. Podrobnejšie o IFP pojednáva samostatná kapitola tejto práce (Podebradský et al., ˇ 2009; Navrátil et al., 2019).
 Fototerapia a mechanoterapia Nízkohladinový laser (LLLT) a ultrazvuk podporujú hojivé procesy v mäkkých tkanivách prostredníctvom fotobiomodulacných a termických ú ˇ cinkov. Tieto modality sú vhodné ˇ pri tendosynovitíde, entezopatii a na podporu hojenia po artikulárnych injekciách (Podeb- ˇ radský et al., 2009).
 1.9.4 Balneoterapia a kúpel’ná liecba ˇ Významnou súcast’ou komplexnej starostlivosti o pacientov s RA je balneoterapia, ˇ ktorá má v našich podmienkach dlhorocnú tradíciu. Ide o využitie prírodných lie ˇ civých vôd, ˇ peloidov (bahna) a klimatických faktorov. Podl’a najnovšej meta-analýzy, ktorú publikovala Aribi et al. (2025), prináša balneoterapia signifikantné zlepšenie v redukcii bolesti a zlepšení kvality života (QoL) u pacientov s reumatickými ochoreniami. Efekt nie je len placebo; termálne minerálne vody pôsobia chemicky (cez kožu sa vstrebávajú minerály ako síra, jód) a fyzikálne (teplo, vztlak). Verhagen et al. (2015) v Cochrane systematickom prehl’ade potvrdzujú, že ``kúpel’ná liecba``(Spa therapy) má pozitívny vplyv na funk ˇ cnú kapacitu pacientov s ˇ RA, pricom benefity pretrvávajú 3 až 6 mesiacov po ukon ˇ cení lie ˇ cebného pobytu. Špecifický ˇ význam majú sírne vody (napr. Piešt’any, Trencianske Teplice). Štúdia autorov Zapolski et ˇ al. (2024) poukázala na to, že sírne kúpele nielen tlmia zápal, ale majú aj pozitívny vplyv na kardiovaskulárne parametre, co je u RA pacientov ( ˇ casto trpiacich aterosklerózou) kl’ú ˇ cové. ˇ Mechanizmus úcinku zah ˇ  ́rna pravdepodobne aj moduláciu oxida ˇ cného stresu. W  ̨ator ( ˇ 2020) uvádzajú, že balneochemické vlastnosti vôd ovplyvnujú redoxný potenciál v tkanivách, ˇ címˇ napomáhajú neutralizácii vol’ných radikálov vznikajúcich pri chronickom zápale.
 1.9.5 Ergoterapia a ochrana k ́lbov Ergoterapia ucí pacienta techniky ``Joint Protection``– ako vykonávat’ bežné ˇ cinnosti ˇ bez zbytocného pret’ažovania k ˇ  ́lbov. Zah ́rna používanie dlahovania na prevenciu deformít, ˇ úpravu domáceho prostredia a používanie kompenzacných pomôcok (napr. špeciálne otvá- ˇ race, zosilnené rú ˇ cky príborov), ktoré šetria drobné k ˇ  ́lby ruky a umožnujú pacientovi zostat’ ˇ nezávislým (Herbert et al., 2022).
 Rehabilitacné ošetrovatel’stvo dop ˇ  ́lna ergoterapeutickú intervenciu o nácvik základ- ˇ ných sebaobslužných aktivít u pacientov s t’ažkými poruchami hybnosti. Ciel’om je maximalizácia funkcnej nezávislosti a prevencia sekundárnych komplikácií z imobility (Kluso ˇ nová ˇ et al., 2014).
 1.9.6 Rehabilitácia a reumatológia – interdisciplinárna spolupráca Optimálny manažment pacienta s RA vyžaduje koordinovanú spoluprácu reumatológa, fyzioterapeuta, ergoterapeuta, rehabilitacného lekára a d’alších špecialistov. Rehabi- ˇ litacná intervencia by mala byt’ integrovaná do lie ˇ cebného plánu od momentu stanovenia ˇ diagnózy a nie až po rozvoji ireverzibilných deformít (Gúth, 2010).
 Vcasná rehabilitácia zah ˇ  ́rna nielen kinezioterapiu a fyzikálnu lie ˇ cbu, ale aj edukáciu ˇ pacienta o ochorení, nácvik kompenzacných stratégií a psychologickú podporu. Pacient by ˇ mal byt’ aktívnym partnerom v liecebnom procese a mat’ realistické o ˇ cakávania ohl’adom ˇ priebehu ochorenia a možností terapie (Gúth, 2010; Špiritovic, ́ 2018).
 1.10 Prognóza reumatoidnej artritídy Prognóza RA je vysoko variabilná a závisí od množstva faktorov. Identifikácia prognostických markerov umožnuje stratifikovat’ pacientov podl’a rizika a prispôsobit’ intenzitu ˇ liecby (Wolfe et al., ˇ 2010).
 1.10.1 Faktory ovplyvnujúce dlhodobý výsledok ˇ Nepriaznivé prognostické faktory zah ́rnajú vysoké titre reumatoidného faktora a anti- ˇ CCP protilátok, prítomnost’ erozívnych zmien v skorom štádiu ochorenia, extraartikulárne manifestácie, vysokú zápalovú aktivitu a funkcnú disabilitu na za ˇ ciatku choroby. Ženy majú ˇ všeobecne horšiu prognózu než muži a genetické faktory (HLA-DR4) sú spojené s agresívnejším priebehom (Wolfe et al., 2010).
 Vcasné zahájenie lie ˇ cby DMARDs, dosiahnutie remisie do 6 mesiacov od za ˇ ciatku ˇ symptómov a striktná kontrola zápalovej aktivity (``Treat to Target``) signifikantne zlepšujú dlhodobé výsledky. Pacienti, ktorí dosiahnu remisiu v ökne príležitosti``, majú dobré šance na dlhodobé funkcné zachovanie bez významného štrukturálneho poškodenia (K. Pavelka, ˇ 2012; McCarthy et al., 2018).
 1.10.2 Mortalita a kvalita života Napriek pokrokom v terapii zostáva stredná d ́lžka života pacientov s RA skrátená v porovnaní s bežnou populáciou, a to približne o 5 až 10 rokov. Hlavnými prícinami smrti sú ˇ kardiovaskulárne ochorenia, infekcie a malignity. Moderná agresívna terapia však viedla k zlepšeniu prežívania v posledných dekádach (Wolfe et al., 2010).
 Kvalita života súvisí predovšetkým s funkcným stavom a kontrolou bolesti. Pacienti ˇ v remisii alebo s nízkou aktivitou ochorenia udávajú kvalitu života porovnatel’nú s bežnou populáciou. Súcasná terapeutická stratégia preto kladie dôraz nielen na kontrolu zápalu, ale ˇ aj na maximalizáciu funkcnej kapacity a aktívne zapojenie pacienta do bežného života (Ole- ˇ járová, 2008).
 